In multi-hop trusted-relay quantum key distribution (QKD), secret key is a
depletable per-link resource, since each hop of a relayed path consumes it. How a
network should behave when that key runs short is unsettled. ITU-T Y.3803 permits
degradation, allowing a symmetric cipher to replace one-time-pad key relay, and
field deployments follow it. A parallel line of security engineering refuses,
treating the same transition as a downgrade attack and failing closed instead.
Neither position has been measured against the other. This paper quantifies the
trade-off. We model a nine-node tactical relay formation with per-link key
reservoirs, under an adversary that raises the quantum bit error rate (QBER) until
affected links abort key generation. We introduce \emph{forced assurance
loss}---the fraction of sessions entitled to an information-theoretically secure
key that are served without one---and the \emph{downgrade-forcing cost}, the
adversary budget required to induce it. Suppressing links equal in number to the
minimum edge cut forces the choice. Fail-closed operation then loses availability,
from 0.812 to 0.304. Fail-degraded operation instead preserves availability at
0.973, but serves 69\% of entitled traffic below its authorised grade. The forcing
cost obeys a two-threshold law: the adversary must both exceed a per-link QBER
margin and hold at least the minimum cut; below either threshold its effect is
negligible. A post-quantum fallback removes the sensitivity of availability to
routing policy, but protects only against quantum-channel attacks. These findings
indicate that the assurance-grade composition of served traffic belongs in QKD
resilience reporting alongside key-supply continuity. The simulator, validated
against four published measurements, is released with all configurations.
